% TOP_PROBLEM: {"id": 20000881, "title": "Tensor growth rate and exact one-dimensional obstructions for constrained differences", "category_id": 2, "category": {"id": 2, "name": "combinatorics", "display_name": "Combinatorics", "description": "Counting problems, graph theory, discrete structures.", "slug": "combinatorics", "order_index": 2, "created_at": "2026-07-31T15:26:25.671Z"}, "status": "partially_solved", "problem_number": "AIM-COMBINATORICS-0006", "set_id": 14}
% FIRST_POSTED: 2026-10-01
% ENTRY_KIND: statement_only
% UnsolvedMath ID 20000881; source catalogue code AIM-COMBINATORICS-0006
% !TeX program = lualatex
% Definition and sources only; this file is not an LLM solution attempt.
\documentclass[11pt,a4paper]{article}
\usepackage[margin=25mm]{geometry}
\usepackage{fontspec}
\setmainfont{lmroman10-regular.otf}[BoldFont=lmroman10-bold.otf,
 ItalicFont=lmroman10-italic.otf,BoldItalicFont=lmroman10-bolditalic.otf]
\usepackage{amsmath,amssymb,mathtools}
\usepackage{unicode-math}
\setmathfont{latinmodern-math.otf}
\AtBeginDocument{\let\setminus\smallsetminus}
\usepackage{xurl,hyperref}
\hypersetup{colorlinks=true,urlcolor=blue}
\setcounter{secnumdepth}{0}
\setlength{\parindent}{0pt}
\setlength{\parskip}{5pt}
\setlength{\emergencystretch}{3em}
\begin{document}
\section*{Tensor growth rate and exact one-dimensional obstructions for constrained differences}\label{20000881}
{\small UnsolvedMath ID 20000881 · AIM-COMBINATORICS-0006 · Combinatorics · AIM Workshop Problem Lists}
\subsection{Definitions and mathematical statement}
Fix a prime $p$ and a subset $D\subseteq\mathbb F_p$ containing $0$.
For an integer $n\ge1$, write $D^n$ for its Cartesian power, considered
inside the vector space $\mathbb F_p^n$. Call $A\subseteq\mathbb F_p^n$
$D$-progression-free if there are no vectors $x\in\mathbb F_p^n$ and
$y\in D^n\setminus\{0\}$ for which
\[
                 x,\quad x+y,\quad x+2y\in A.
\]
All coordinates and scalar multiples are computed in $\mathbb F_p$.
The restriction is on the difference vector: different coordinates of
$y$ may be different elements of $D$, and some may vanish.
Define the extremal size
\[
 a_D(n)=\max\{|A|:A\subseteq\mathbb F_p^n\text{ is }D\text{-progression-free}\}.
\]
The problem asks for a characterization of the pairs $(p,D)$ for which
there exists a real constant $C=C(p,D)<p$ such that
$a_D(n)\le C^n$ for every positive integer $n$. The same $C$ must work
in all dimensions. Thus a density estimate $a_D(n)/p^n\to0$ alone
does not answer the question: the requested saving is exponential in $n$.
The original algebraic formulation is retained also in characteristic two,
where $x+2y=x$; no extra requirement that the three displayed points be
distinct is imposed. For odd $p$, their distinctness follows from $y\ne0$.
The set $D=\{0\}$ is included as a boundary case in the classification.
\subsection{Short English statement}
Which allowed coordinate differences force every progression-free subset of a finite vector space to have an exponential loss in size?
\subsection{Sources}
\begin{itemize}
\item[S1] ULAM.AI and the UnsolvedMath Contributors, supplied problems.json, record 20000881 (AIM-COMBINATORICS-0006), AIM Workshop Problem Lists. Catalogue identity and source transcription; CC BY 4.0.
\url{https://huggingface.co/datasets/ulamai/UnsolvedMath}.
\item[S2] American Institute of Mathematics, High-dimensional phenomena in discrete analysis, item 2.3. Original workshop question underlying AIM-COMBINATORICS-0006.
\url{http://aimpl.org/highdimdiscrete/2/}.
\end{itemize}
\end{document}
